\section{Model}\label{sec:model}

A stylized model captures what the set-aside does and, as important, what it does not do. $N$ potential bidders compete for a single indivisible contract. Each bidder $i$ has type $k_i \in \{\text{SME}, \neg\}$ and draws a private cost of supply $c_i$ from a type-specific distribution $F_c^k$, with $F_c^{\text{SME}}$ and $F_c^{\neg}$ differing in ways the data are allowed to determine.

The auction format is the identifying payload. BEC implements procurement through a Preg\~ao---a descending-clock electronic English-reverse auction, distinct both from first-price sealed-bid designs (Convite, in the smaller-value Brazilian segment) and from the Dutch descending-price format more commonly studied in the auction literature. The clock opens at the buyer's reference price $p^{\mathrm{ref}}$, firms post downward-revisable offers until only one firm remains, and that firm supplies the contract at the clearing price. What makes this format analytically convenient is a standard result of \citet{haile2003}: under independent private values, the weakly dominant strategy is to exit the auction when the running clock price reaches one's own cost. Drop-out prices of losers therefore \emph{point-identify} their costs, without requiring the researcher to invert a bid function or assume a parametric form for the bidding strategy. The winner, who never drops out before the clock stops, leaves only an upper bound $c_{\text{win}} \le c_{(2)}$ (the cost of the second-to-last firm to exit)---and that upper bound is exactly what the Turnbull robustness check in Section~\ref{sec:robustness} uses to deliver a non-parametric MLE of $F_c^k$. This point-identification under Preg\~ao is the reason a structural analysis is available here where the cross-sectional first-price-sealed-bid settings of \citet{marion2007}, \citet{krasnokutskaya2011}, and \citet{atheyseira2011} have relied on cross-equation constraints to disentangle channels.

\paragraph{Why this framework.} The choice of asymmetric IPV with the Krasnokutskaya UH correction and Athey--Seira reduced-form entry is dictated by the question and by what the BEC-Preg\~ao data identify. Three alternatives merit explicit comment. \emph{Common-value or affiliated-private-value formulations} \citep{milgrom1982, hong2003} are most natural where bidders share an unknown component of the contract's value---e.g., construction contracts with incompletely observed specifications. The CADMAT-class restriction strips cross-item heterogeneity at the catalog level, and the cross-modality consistency check tolerates within-auction cost correlation up to $\rho_c \le \ipvRhoMax$ without rejecting the IPV restriction (Section~\ref{sec:robustness}). \emph{Nested entry-bidding equilibria} \`a la \citet{krasnokutskayaseim2011} recover entry costs structurally from a free-entry indifference condition; the deliberately weaker reduced-form entry adopted here---Pre/Post bidder counts as observed equilibrium primitives in the \citet{atheyseira2011} spirit---keeps the welfare statement anchored to the legal trigger rather than to a participation-cost identity the trigger does not directly identify. \emph{Open-format or FPSB equivalents} are bounded explicitly under \citet{maskinriley2000} in Section~\ref{sec:fpsb_bound}, with the V3-vs-V0 ranking preserved at the upper bound. The framework is therefore the lightest specification that sustains the within-auction-vs-entry decomposition while leaving alternative readings to the robustness battery.

The model relies on three substantive restrictions, stated formally below and defended in turn with the diagnostic evidence in Section~\ref{sec:identification}.

\begin{assumption}[Independent Private Values, asymmetric across types]\label{a:ipv}
Costs $c_i$ are drawn independently across bidders. Conditional on type $k$, $c_i \sim F_c^k$ with density $f_c^k$ on support $[\underline{c}, \overline{c}] \subset \mathbb{R}_+$. The SME and non-SME distributions are not constrained to coincide.
\end{assumption}

Assumption~\ref{a:ipv} permits but does not impose stochastic dominance between types; the empirical median-cost gap of \medCostGapLo--\medCostGapHi{} of the reference price across (pharma $\times$ type) cells (Section~\ref{sec:identification}) supports asymmetry directly, and symmetric IPV would require this gap to vanish.

\begin{assumption}[Multiplicative auction-level heterogeneity]\label{a:uh}
The normalized log-bid decomposes as $y_{it} = \log(b_{it} / p_t^{\mathrm{ref}}) = a_t + e_{it}$, where $a_t$ is an auction-level scale common to all bidders in auction $t$, and $e_{it}$ is bidder-specific and independent across bidders within $t$. Both $a_t$ and $e_{it}$ are mean-zero with finite variances $\sigma_a^{2,k}$ and $\sigma_e^{2,k}$ that may vary by stratum.
\end{assumption}

Assumption~\ref{a:uh} is the \citet{krasnokutskaya2011} restriction separating a common scale from bidder-specific private values; multiplicativity is the natural form when item complexity or ANVISA burden scales costs uniformly across firms. The Turnbull NPMLE of Section~\ref{sec:robustness} dispenses with Assumption~\ref{a:uh} altogether and corroborates the baseline within \turnbullDeviationPct~percent of the total effect.

\begin{assumption}[Set-aside admissibility; SME cost distribution allowed to differ pre/post]\label{a:setaside}
Under the open regime both SMEs and non-SMEs enter with strictly positive probability. Under the SME-only regime, only SMEs are admissible. Within a short window around the cutoff, the underlying technology of supply is stable, so the non-SME cost distribution $F_c^{\neg}$ is constant across periods. The SME cost distribution is allowed to differ between regimes: write $F_c^{\text{SME},\tau}$ for $\tau \in \{\text{Pre},\text{Post}\}$. We do not model the mechanism that generates this difference (a formal entry-selection equilibrium in the spirit of \citet{krasnokutskayaseim2011} is left for future work); the empirical content of $F_c^{\text{SME,Post}} \neq F_c^{\text{SME,Pre}}$ is treated as an upper bound on the combined contribution of selection, technology shifts within the SME segment, and any other mechanism that would change the empirical cost distribution of active SME bidders between the two regimes.
\end{assumption}

\begin{proposition}[Identification of $F_c^k$ under ascending-clock IPV]\label{prop:identification}
Under Assumptions~\ref{a:ipv}--\ref{a:setaside}, drop-out exit prices in the Preg\~ao English-reverse format point-identify the type-specific cost distribution $F_c^k$ \citep{haile2003}, up to the auction-level scale $a_t$ recovered by the variance decomposition in Assumption~\ref{a:uh} \citep{krasnokutskaya2011}. The post-policy SME distribution $F_c^{\text{SME,Post}}$ is point-identified separately from $F_c^{\text{SME,Pre}}$ from post-period drop-out exits; strict invariance ($F_c^{\text{SME,Post}} = F_c^{\text{SME,Pre}}$) is recovered as a nested benchmark by imposing the equality.
\end{proposition}

Assumption~\ref{a:setaside} is the structural counterpart of the parallel-trends assumption in reduced-form DiD, applied to the SME segment only: the non-SME cost distribution is required to be invariant across periods (a testable restriction, examined in Section~\ref{sec:identification}), while the SME cost distribution is allowed to differ. The main specification estimates $F_c^{\text{SME,Pre}}$ and $F_c^{\text{SME,Post}}$ separately from each period's drop-out exits and reports the difference; the strict-invariance benchmark imposes equality and re-runs the same simulation. Neither specification tests a formal selection model---we estimate distributions, we do not derive equilibrium selection from a free-entry primitive in the \citet{krasnokutskayaseim2011} sense. The empirical difference $\hat F_c^{\text{SME,Post}} - \hat F_c^{\text{SME,Pre}}$ is therefore a measurement upper bound on the policy effect on the active SME cost distribution; the gap between the two welfare readings (main vs.\ strict invariance) is the corresponding bound on how much the welfare conclusion depends on attributing this empirical difference to selection rather than to no real change in the underlying object. The primitive-invariance KS test on non-SME $F_c$ (Table~\ref{tab:v3_primitive_invariance}) rejects strict invariance in pharma Preg\~ao ($D = \ksPharmaPregao$) and non-pharma Convite ($D = \ksNpConvite$) at the operational threshold $D < \ksDistThresh$, and fails to reject in the other two strata. The cross-modality check (Table~\ref{tab:v3_cross_modality}) provides an additional consistency test.

\paragraph{Two readings of $F_c^{\text{SME,Post}}$ as bounds.} The framework delivers two readings of the policy effect, which we report as bounds rather than as a tested specification. The \emph{main} reading uses $\hat F_c^{\text{SME,Post}}$ as estimated from post-period exits, which absorbs into the SME cost primitive whatever empirical difference exists between SME bidders observed Pre and Post. The \emph{strict-invariance} reading imposes $F_c^{\text{SME,Post}} = F_c^{\text{SME,Pre}}$ and attributes the entire post-cutoff price movement to changes in bidder counts. Neither reading is a model-tested object: the main reading does not derive the empirical $\hat F_c^{\text{SME,Post}} - \hat F_c^{\text{SME,Pre}}$ gap from a free-entry equilibrium, and the strict-invariance reading imposes equality on a quantity the KS tests in Tables~\ref{tab:v3_primitive_invariance} and~\ref{tab:v3_uh_invariance} reject in three of the four Preg\~ao strata after UH cleaning (KS distances of \ksNpPregaoUH{} non-pharma and \ksPharmaPregaoUH{} pharma). The pair therefore brackets the welfare statement under two contrasting attributions of the same empirical fact---the post-policy SME distribution being measurably different from the pre-policy one. The non-pharmaceutical welfare ranking holds under both bounds; the pharmaceutical welfare ranking holds under main but reverses under strict invariance, and the gap between the two readings is itself the substantive object Section~\ref{sec:robustness} reports.

\paragraph{Equilibrium, entry, and counterfactuals.} The auction format is interpreted as ascending-clock IPV: under Assumption~\ref{a:ipv}, exit-at-cost is weakly dominant, the surviving bidder pays the second-lowest exit, and the expected transaction price equals $\mathbb{E}[c_{(2)}]$. Revenue equivalence \citep{vickrey1961, maskinriley2000} extends the same object to FPSB; \citet{krasnokutskaya2011} provides the analog argument used here. The Bayes--Nash counterfactual reduces to Monte Carlo simulation over draws from the estimated $F_c^k$: stratum-specific bidder counts from the type-specific Poisson arrival process, costs from $F_c^k$, sorted to read off $c_{(1)}$ (production cost) and $c_{(2)}$ (expected transaction price). Equilibrium arrival rates enter as primitives, not derived from free entry: observed Pre counts $(N^{\text{SME}}_{\text{Pre}}, N^{\neg}_{\text{Pre}})$ are status-quo arrivals; Post counts $(N^{\text{SME}}_{\text{Post}}, 0)$ are equilibrium arrival under the restricted regime, in the \citet{atheyseira2011} spirit. The implied zero-profit entry costs (Table~\ref{tab:v3_entry_cost}, \entryCostMin--\entryCostMax{} per bidder) are per-auction calibration margins, not jointly identified fixed costs of the kind in \citet{krasnokutskayaseim2011}; the decomposition that follows is counterfactual-simulation accounting under the maintained assumptions, and the welfare-cost expression of Section~\ref{sec:welfare} prices the simulated price effect at the MCPF $\lambda$.\footnote{Independence-across-auctions is assumed at the level of cost draws conditional on type and stratum. Within-auction correlation enters through the auction-level scale $a_t$ (Assumption~\ref{a:uh}); the cluster bootstrap of Section~\ref{sec:results} resamples at the auction level to preserve this correlation in the standard errors.}

\paragraph{The decomposition.} Three counterfactual scenarios anchor the exercise:
\begin{itemize}
\item $S_1$: open regime at the pre-period pool. The baseline.
\item $S_2$: SME-only at the pre-period SME pool. The non-SME type is removed; the SME pool is held at its observed pre-policy arrival rate. Within-auction component isolated.
\item $S_3$: SME-only at the post-period SME pool. Uses the observed post-policy arrival rate and the post-policy SME cost distribution. Full policy effect with endogenous entry.
\end{itemize}
The simulated total effect on the second-order statistic decomposes arithmetically as
\begin{equation}
\underbrace{p_{S_3} - p_{S_1}}_{\text{simulated total effect}}
\;=\; \underbrace{(p_{S_2} - p_{S_1})}_{\text{within-auction component (sheltered bidding)}}
   \,+\, \underbrace{(p_{S_3} - p_{S_2})}_{\text{participation-margin component}}.
\end{equation}
Both components are properties of the simulation given $F_c^k$ and the observed entry pattern; they sum to the simulated total effect by construction. The within-auction component can exceed the realized total effect when the participation-margin component partially offsets it, which is what the data deliver in non-pharmaceutical markets.

\paragraph{What is new here.} The within-auction component is implicit in the asymmetric-IPV structural literature on bidder preferences: \citet{krasnokutskaya2011} and \citet{krasnokutskayaseim2011} recover it on FPSB highway auctions via cross-equation constraints on the bid function, and \citet{athey2013} computes the analog object under set-asides and subsidies in U.S.\ timber FPSB. Two features distinguish the present treatment. First, the descending-clock format point-identifies the cost distribution from exit prices alone \citep{haile2003}, so the within-auction component does not require inverting a bid function nor a cross-equation restriction---it is delivered by the auction format itself. Second, ascending-clock IPV means exit-at-cost remains weakly dominant regardless of pool composition, so the within-auction shift produced by removing a bidder type involves no behavioral re-adjustment by surviving bidders; this is mechanically different from FPSB, where bidders re-bid against a thinner distribution and the analog object combines an order-statistic shift with a strategic markup response that the format does not separately identify. The definition that follows formalizes the object so that the characterization survives outside the BEC-SP application; the post-definition paragraph develops the identification under ascending-clock IPV explicitly.

\begin{definition}[Sheltered bidding]\label{def:sheltered}
Consider an asymmetric IPV auction in which the buyer restricts admissibility to bidder type $k$. Let $p_{\text{open}}$ denote the expected transaction price under the open regime with both types entering at observed equilibrium arrival rates, and let $p_{\text{shelter}}$ denote the expected transaction price under the type-$k$-only regime at the same on-type pool size as the open regime. The \emph{sheltered-bidding component} of the policy effect is the equilibrium price shift
\[
\Delta_{\text{shelter}} \;\equiv\; p_{\text{shelter}} - p_{\text{open}},
\]
identified separately from the \emph{participation-margin component} $\Delta_{\text{entry}} \equiv p_{\text{endo}} - p_{\text{shelter}}$ that captures the equilibrium response of type-$k$ entry to the restriction. In the BEC-SP notation, $\Delta_{\text{shelter}} = p_{S_2} - p_{S_1}$ and $\Delta_{\text{entry}} = p_{S_3} - p_{S_2}$.
\end{definition}

\paragraph{Identification under ascending-clock IPV.} Under Assumptions~\ref{a:ipv}--\ref{a:setaside}, $\Delta_{\text{shelter}}$ is point-identified from drop-out exits via Proposition~\ref{prop:identification}; the expected second-order statistic shifts as a function of which order statistic is realized at each pool composition, with no active behavioral adjustment by surviving bidders. By contrast, the FPSB analog combines an order-statistic shift with a strategic re-bid against a thinner distribution---the bid function $\beta(c; F^k_c, N)$ depends on the entire competitor type distribution, so removing a type shifts the symmetric-equilibrium bid even before any order-statistic effect---requiring the cross-equation machinery of \citet{krasnokutskaya2011} or the strategic-equilibrium recovery of \citet{lebrun1999} to separate the two channels. The descending-clock format delivers the decomposition mechanically.

The component is structural in the sense that any policy intervention that removes a bidder type while holding the surviving pool fixed will produce a $\Delta_{\text{shelter}}$ predictable from $F_c^k$ and $F_c^{\neg}$ alone, under the auction format's exit-at-cost incentive---procurement set-asides, exclusionary clearance rules, supplier disqualification regimes, or any restriction that operates on the admissibility margin within an ascending-clock environment. The FPSB upper bound of \citet{maskinriley2000} (Section~\ref{sec:fpsb_bound}) checks how much the result would move under the alternative format.

\paragraph{When sheltered bidding dominates.} The decomposition predicts $|\Delta_{\text{shelter}}| > |\Delta_{\text{entry}}|$ when two conditions hold jointly: the off-type pool was thick enough in the open regime to materially shift the second-order statistic ($F_c^{\neg}$ contributes with non-negligible mass to $\mathbb{E}[c_{(2)}]$), and the type-$k$ entry response to the restriction is bounded relative to the off-type removal effect. Both conditions are satisfied in the thick standardized non-pharmaceutical case of the BEC-SP data, where $\Delta_{\text{shelter}}/\Delta_{\text{total}} \approx \bneShareIntensNp$~percent. The pharmaceutical case sits closer to the boundary on both margins: a thinner off-type pool reduces $|\Delta_{\text{shelter}}|$, and the empirical pre-vs-post difference in $\hat F_c^{\text{SME}}$ introduces an additional source of variation in the entry margin that the static framework does not bound under any specific mechanism. The bifurcation between the model-stable non-pharmaceutical decomposition and the model-sensitive pharmaceutical one is therefore a prediction of the framework, not an artifact of the application, and the framework provides the diagnostic policy-makers need when choosing between bidder-exclusion and price-preference instruments across heterogeneous procurement categories.
